\documentclass[UTF8, 12pt, fontset=fandol, oneside]{ctexart}
\usepackage{researchnotes}

\title{强化学习：导论}
\author{}
\date{}

\begin{document}

\maketitle

\section*{第 8 章\quad 表格型方法中的规划与学习}

本章中，我们建立一种统一观点，来看待那些需要环境模型的强化学习方法，例如动态规划和启发式搜索；也来看待那些不需要模型即可使用的方法，例如蒙特卡洛方法和时序差分方法。
这两类方法分别称为\textbf{基于模型}的强化学习方法和\textbf{无模型}的强化学习方法。
基于模型的方法主要依赖\textbf{规划}，而无模型方法主要依赖\textbf{学习}。
虽然这两类方法之间确实存在真实差别，但它们也有很大的相似性。
特别地，两类方法的核心都是价值函数的计算。
此外，所有这些方法都基于向前看未来事件、计算一个备份值，再把它用作近似价值函数的更新目标。
本书前面曾把蒙特卡洛方法和时序差分方法作为不同的替代方案来介绍，然后又说明它们可以由 $n$ 步方法统一起来。
本章的目标，是对基于模型方法和无模型方法作类似的整合。
既然前面已经把它们确立为不同方法，现在我们就来考察它们能够在多大程度上彼此混合。

\subsection*{8.1\quad 模型与规划}

所谓环境的\textbf{模型}，是指智能体可以用来预测环境将如何响应其动作的任何东西。
给定一个状态和一个动作，模型会产生对随后的下一个状态与下一个奖励的预测。
如果模型是随机性的，那么就会有若干可能的下一个状态和下一个奖励，每一种都有某个发生概率。
有些模型会给出所有可能性及其概率的描述；我们称这样的模型为\textbf{分布模型}。
另一些模型只按照这些概率抽样地产生其中一种可能性；我们称这样的模型为\textbf{样本模型}。
例如，考虑对十二个骰子的点数和建模。
分布模型会产生所有可能的点数和以及它们发生的概率，而样本模型会按照这个概率分布产生一个单独抽出的点数和。
动态规划中所假设的那类模型，也就是对 MDP 动态 $p(s',r|s,a)$ 的估计，是分布模型。
第 5 章二十一点例子中使用的模型则是样本模型。
分布模型比样本模型更强，因为它们总能用来产生样本。
然而，在许多应用中，得到样本模型比分布模型容易得多。
十二个骰子就是一个简单例子。
写一个计算机程序来模拟掷骰并返回点数和很容易，而列出所有可能点数和及其概率则更困难，也更容易出错。

模型可以用来模仿或模拟经验。
给定一个起始状态和动作，样本模型会产生一个可能的转移，而分布模型会生成所有可能转移，并按它们发生的概率加权。
给定一个起始状态和一个策略，样本模型可以产生一整个幕，而分布模型可以生成所有可能的幕及其概率。
在任一情形下，我们都说这个模型被用来\textbf{模拟}环境并产生\textbf{模拟经验}。

“规划”这个词在不同领域中有若干不同用法。
我们用它指任何以模型为输入、并产生或改进与被建模环境交互的策略的计算过程：
\[
\text{模型}\quad \xrightarrow{\quad \text{规划}\quad}\quad \text{策略}.
\]

在人工智能中，按我们的定义，有两种不同的规划途径。
\textbf{状态空间规划}包括本书采用的方法，它主要被看作是在状态空间中搜索最优策略或通往目标的最优路径。
动作导致从一个状态到另一个状态的转移，而价值函数则在状态上计算。
在我们称为\textbf{计划空间规划}的方法中，规划反而是在计划空间中搜索。
算子把一个计划变换成另一个计划，而如果定义价值函数，它也是定义在计划空间上的。
计划空间规划包括演化方法和“部分序规划”；后者是人工智能中常见的一类规划，其中步骤顺序在规划的所有阶段都不完全确定。
计划空间方法难以高效应用于强化学习所关注的随机序贯决策问题，因此我们不再进一步考虑它们（但可参见 Russell and Norvig, 2010）。

本章给出的统一观点是：所有状态空间规划方法都共享一个共同结构，而这个结构也出现在本书所介绍的学习方法中。
展开这个观点需要本章其余内容，但其中有两个基本思想：
（1）所有状态空间规划方法都会把计算价值函数作为改进策略的关键中间步骤；
（2）它们通过对模拟经验施加更新或备份操作来计算价值函数。
这种共同结构可以图示为：
\[
\text{模型}
\longrightarrow
\text{模拟经验}
\xrightarrow{\quad \text{备份}\quad}
\text{价值}
\longrightarrow
\text{策略}.
\]

动态规划方法显然符合这一结构：它们扫过状态空间，为每个状态生成可能转移的分布。
然后用每个分布计算备份后的价值（更新目标），并更新该状态的估计价值。
本章我们将说明，各种其他状态空间规划方法也符合这一结构；不同的具体方法只是在所做更新的种类、更新的顺序，以及备份信息保留多久这些方面有所差异。

以这种方式看待规划方法，会突出它们与本书已经描述过的学习方法之间的关系。
学习和规划方法的核心，都是通过备份式更新操作来估计价值函数。
差别在于，规划使用模型生成的模拟经验，而学习方法使用环境生成的真实经验。
当然，这个差别还会导致若干其他差异，例如性能如何评价，以及经验能够以多灵活的方式生成。
但共同结构意味着，许多思想和算法可以在规划与学习之间转移。
特别地，在许多情况下，一个学习算法可以替代规划方法中的关键更新步骤。
学习方法只需要经验作为输入；在许多情况下，它们既可以用于模拟经验，也可以用于真实经验。
下面的框给出了一个基于一步表格型 Q 学习、并从样本模型中随机抽样的简单规划方法。
我们称这个方法为\textbf{随机样本一步表格型 Q 规划}；在与一步表格型 Q 学习对真实环境收敛到最优策略相同的条件下，它也会收敛到模型的最优策略，也就是说，在步骤 1 中每个状态--动作对必须被选择无穷多次，并且 $\alpha$ 必须随时间适当减小。

\begin{center}
\fbox{\begin{minipage}{0.88\textwidth}
\noindent\textbf{随机样本一步表格型 Q 规划}

\vspace{0.4em}
\noindent 永远循环：
\begin{enumerate}
\item 随机选择一个状态 $S\in\mathcal{S}$ 和一个动作 $A\in\mathcal{A}(S)$；
\item 把 $S,A$ 送入样本模型，并得到一个样本下一个奖励 $R$ 与一个样本下一个状态 $S'$；
\item 对 $S,A,R,S'$ 应用一步表格型 Q 学习：
\[
Q(S,A)\leftarrow Q(S,A)+\alpha\left[R+\gamma\max_a Q(S',a)-Q(S,A)\right].
\]
\end{enumerate}
\end{minipage}}
\end{center}

除了规划方法与学习方法的统一观点之外，本章第二个主题是以小的、增量的步骤进行规划的好处。
这样一来，规划可以随时被中断或重新定向，浪费的计算很少；这似乎是把规划与行动以及模型学习高效混合起来的关键要求。
如果问题大到无法精确求解，那么即使在纯规划问题上，非常小步的规划也可能是最高效的途径。

\subsection*{8.2\quad Dyna：规划、行动与学习的整合}

当规划在与环境交互的同时在线进行时，会出现许多有趣问题。
从交互中获得的新信息可能改变模型，并因此与规划发生相互作用。
也许需要以某种方式把规划过程定制到当前正在考虑、或预计不久后将考虑的状态或决策上。
如果决策制定和模型学习都是计算密集型过程，那么可用的计算资源可能需要在二者之间分配。
为了开始探索这些问题，本节介绍 Dyna-Q：一种简单架构，它整合了在线规划智能体所需的主要功能。
每个功能都以简单、几乎平凡的形式出现在 Dyna-Q 中。
在后续各节中，我们会展开实现每个功能的若干替代方式，以及它们之间的权衡。
目前，我们只想说明这些思想并激发直觉。

在规划智能体内部，真实经验至少有两个作用：它可以用来改进模型，使模型更准确地匹配真实环境；也可以直接用前几章讨论过的强化学习方法来改进价值函数和策略。
前者称为\textbf{模型学习}，后者称为\textbf{直接强化学习}（direct RL）。
经验、模型、价值与策略之间可能的关系，可以概括为右侧那类图示。
每条箭头都表示一种影响和推定的改进关系。
注意，经验既可以直接改进价值函数和策略，也可以通过模型间接改进它们。
后者有时称为\textbf{间接强化学习}，它就是规划所涉及的内容。

直接方法和间接方法各有优缺点。
间接方法通常更充分地利用有限经验，因此用更少环境交互就能得到更好的策略。
另一方面，直接方法简单得多，而且不受模型设计偏差的影响。
有人认为间接方法总是优于直接方法，也有人认为直接方法解释了大多数人类和动物学习。
心理学和人工智能中的相关争论，还涉及认知相对于试错学习的重要性，以及深思熟虑式规划相对于反应式决策的重要性（从心理学角度看这些问题，可见第 14 章）。
我们的观点是，这些争论中各替代方案之间的对比被夸大了；认识到两边的相似性，比把它们对立起来能带来更多洞见。
例如，本书一直强调动态规划和时序差分方法之间的深层相似性，尽管前者是为规划设计的，后者是为无模型学习设计的。

Dyna-Q 包含上图中显示的所有过程：规划、行动、模型学习和直接强化学习，并且它们持续发生。
规划方法就是第 161 页的随机样本一步表格型 Q 规划方法。
直接强化学习方法是一步表格型 Q 学习。
模型学习方法也基于表格，并假定环境是确定性的。
在每次转移
\[
S_t,A_t\to R_{t+1},S_{t+1}
\]
之后，模型在 $S_t,A_t$ 对应的表项中记录这样一个预测：$R_{t+1},S_{t+1}$ 将确定地随之出现。
因此，如果用以前经历过的状态--动作对查询模型，它就直接返回最后一次观察到的下一个状态和下一个奖励，作为它的预测。
在规划期间，Q 规划算法只从以前经历过的状态--动作对中随机抽样（步骤 1），所以模型从不会被问到它没有信息的状态--动作对。

Dyna 智能体的总体架构如图 8.1 所示，Dyna-Q 算法是其中一个例子。
中间一列表示智能体与环境之间的基本交互，它产生一条真实经验轨迹。
图左侧箭头表示直接强化学习，它在真实经验上操作，以改进价值函数和策略。
右侧是基于模型的过程。
模型由真实经验学习而来，并产生模拟经验。
我们用\textbf{搜索控制}一词指这样一个过程：为模型生成的模拟经验选择起始状态和动作。
最后，规划通过把强化学习方法应用到模拟经验上来完成，就好像这些模拟经验真的发生过一样。
通常，如 Dyna-Q 中那样，同一个强化学习方法既用于从真实经验中学习，也用于从模拟经验中规划。
因此，强化学习方法是学习和规划二者的“最终共同通路”。
学习与规划深度整合，因为它们共享几乎全部相同机制，唯一差别只是经验来源不同。

\begin{figure}[htbp]
\centering
\includegraphics[width=0.72\textwidth]{tex/强化学习/figures/math-9/figure-8-1.png}
\caption{Dyna 的一般架构。真实经验在环境与策略之间来回传递，它以很类似于环境模型所生成模拟经验的方式，影响策略与价值函数。}
\end{figure}

从概念上说，在 Dyna 智能体中，规划、行动、模型学习和直接强化学习同时并行发生。
不过，为了在串行计算机上具体实现，我们要完全规定它们在一个时间步内发生的顺序。
在 Dyna-Q 中，行动、模型学习和直接强化学习过程只需要少量计算，我们假定它们只消耗时间的一小部分。
每一步剩余的时间可以用于规划过程，而规划本质上是计算密集型的。
假设在每一步中，在行动、模型学习和直接强化学习之后，还有时间完成 Q 规划算法的 $n$ 次迭代（步骤 1--3）。
下面框中的 Dyna-Q 伪代码里，$Model(s,a)$ 表示状态--动作对 $(s,a)$ 的表项内容，也就是预测的下一个状态和奖励。
直接强化学习、模型学习和规划分别由步骤 (d)、(e) 和 (f) 实现。
如果省略 (e) 和 (f)，剩下的算法就是一步表格型 Q 学习。

\begin{center}
\fbox{\begin{minipage}{0.88\textwidth}
\noindent\textbf{表格型 Dyna-Q}

\vspace{0.4em}
\noindent 对所有 $s\in\mathcal{S}$ 和 $a\in\mathcal{A}(s)$ 初始化 $Q(s,a)$ 与 $Model(s,a)$。

\noindent 永远循环：

\quad (a) $S\leftarrow$ 当前（非终止）状态；

\quad (b) $A\leftarrow \varepsilon$-$greedy(S,Q)$；

\quad (c) 采取动作 $A$；观察得到的奖励 $R$ 和状态 $S'$；

\quad (d) $Q(S,A)\leftarrow Q(S,A)+\alpha\left[R+\gamma\max_a Q(S',a)-Q(S,A)\right]$；

\quad (e) $Model(S,A)\leftarrow R,S'$（假定环境是确定性的）；

\quad (f) 重复 $n$ 次：

\qquad $S\leftarrow$ 随机的、以前观察过的状态；

\qquad $A\leftarrow$ 以前在 $S$ 中采取过的随机动作；

\qquad $R,S'\leftarrow Model(S,A)$；

\qquad $Q(S,A)\leftarrow Q(S,A)+\alpha\left[R+\gamma\max_a Q(S',a)-Q(S,A)\right]$。
\end{minipage}}
\end{center}

\noindent\textbf{例 8.1：Dyna 迷宫}
考虑图 8.2 中嵌入小图所示的简单迷宫。
在 47 个状态中的每一个状态下，都有上、下、右、左四个动作；这些动作会确定性地把智能体带到相应的相邻状态，除非移动被障碍物或迷宫边界挡住，在这种情况下智能体保持在原处。
所有转移上的奖励都是 0，只有进入目标状态的转移奖励为 $+1$。
到达目标状态（G）后，智能体回到起始状态（S）开始新的一幕。
这是一个折扣的分幕任务，其中 $\gamma=0.95$。

图 8.2 的主体部分给出了一个实验中的平均学习曲线；在这个实验中，Dyna-Q 智能体被应用于该迷宫任务。
初始动作价值为 0，步长参数为 $\alpha=0.1$，探索参数为 $\varepsilon=0.1$。
当在动作之间贪心选择且出现平局时，随机打破平局。
这些智能体的差别在于每个真实步之后执行的规划步数 $n$。
对每个 $n$，曲线显示智能体在每一幕中到达目标所用的步数，并在实验的 30 次重复上取平均。
在每次重复中，随机数生成器的初始种子在各算法之间保持相同。
因此，对所有 $n$ 的取值，第一幕完全相同（大约 1700 步），其数据没有显示在图中。
第一幕之后，所有 $n$ 的表现都得到改进，但 $n$ 越大，改进越快。
回忆 $n=0$ 的智能体是不规划的智能体，只使用直接强化学习（一步表格型 Q 学习）。
在这个问题上，它是迄今最慢的智能体，尽管参数值（$\alpha$ 和 $\varepsilon$）已经为它优化过。
不规划的智能体大约用了 25 幕达到 $\varepsilon$-最优表现，而 $n=5$ 的智能体大约用了 5 幕，$n=50$ 的智能体只用了 3 幕。

\begin{figure}[htbp]
\centering
\includegraphics[width=0.78\textwidth]{tex/强化学习/figures/math-9/figure-8-2.png}
\caption{一个简单迷宫（嵌入小图），以及每个真实步中规划步数 $n$ 不同的 Dyna-Q 智能体的平均学习曲线。任务是尽可能快地从 S 走到 G。}
\end{figure}

图 8.3 说明了为什么规划智能体比不规划智能体更快找到解。
图中显示的是 $n=0$ 和 $n=50$ 的智能体在第二幕进行到一半时找到的策略。
没有规划（$n=0$）时，每一幕只会给策略增加一步，因此到目前为止只学到一步，也就是最后一步。
有规划时，第一幕期间同样只学到一步，但在第二幕中已经发展出一个广泛的策略；到第二幕结束时，这个策略几乎会一直回溯到起始状态。
这个策略是在智能体仍在起始状态附近徘徊时，由规划过程建立起来的。
到第三幕结束时，完整的最优策略已经被找到，并且达到完美表现。

\begin{figure}[htbp]
\centering
\includegraphics[width=0.76\textwidth]{tex/强化学习/figures/math-9/figure-8-3.png}
\caption{在第二幕中途，规划型与非规划型 Dyna-Q 智能体找到的策略。箭头表示每个状态中的贪心动作；若某个状态没有显示箭头，则表示它的所有动作价值相等。黑色方块表示智能体的位置。}
\end{figure}

在 Dyna-Q 中，学习和规划由完全相同的算法完成；用于学习时它作用于真实经验，用于规划时它作用于模拟经验。
由于规划以增量方式进行，把规划和行动混合起来是很直接的。
二者都尽可能快地进行。
智能体总是具有反应性，也总是在深思熟虑：它会立即响应最新的感知信息，同时也总是在后台进行规划。
模型学习过程同样在后台持续进行。
随着新信息被获得，模型被更新得更接近现实。
当模型改变时，正在进行的规划过程会逐渐计算出一种不同的行为方式，以匹配新模型。

\noindent\textbf{练习 8.1}
图 8.3 中的不规划方法看起来特别差，因为它是一步方法；使用多步自举的方法会表现得更好。
你认为第 7 章中的某个多步自举方法能像 Dyna 方法一样好吗？
请解释理由。

\subsection*{8.3\quad 当模型是错的}

在上一节的迷宫例子中，模型中的变化相对温和。
模型一开始是空的，然后只填入完全正确的信息。
一般来说，我们不能指望总是如此幸运。
模型可能是不正确的，因为环境是随机的，而我们只观察到了有限数量的样本；也可能因为模型是用会泛化得不完美的函数近似来学习的；还可能仅仅因为环境已经发生变化，而它的新行为尚未被观察到。
当模型不正确时，规划过程很可能计算出一个次优策略。

在某些情况下，规划计算出的次优策略会很快导致建模错误被发现并修正。
当模型在某种意义上是乐观的，也就是预测出比实际可能更大的奖励或更好的状态转移时，这种情况往往会发生。
规划得到的策略会试图利用这些机会，并在这样做时发现它们并不存在。

\noindent\textbf{例 8.2：阻塞迷宫}
图 8.4 给出了一个迷宫例子，用来说明这种相对较轻的建模错误以及从中恢复的情况。
最初，如图左上所示，从起点到目标有一条在障碍右侧的短路径。
1000 个时间步之后，如图右上所示，这条短路径被“阻塞”，而沿障碍左侧打开了一条更长的路径。
图中的曲线显示了一个 Dyna-Q 智能体和一个稍后描述的增强型 Dyna-Q+ 智能体的平均累积奖励。
曲线的第一部分显示，两个 Dyna 智能体都在 1000 步之内找到了短路径。
当环境改变时，曲线变平，表示在一段时间内智能体没有得到奖励，因为它们在障碍后面徘徊。
然而，过一段时间后，它们能够找到新的开口和新的最优行为。

\begin{figure}[htbp]
\centering
\includegraphics[width=0.72\textwidth]{tex/强化学习/figures/math-9/figure-8-4.png}
\caption{Dyna 智能体在阻塞任务上的平均表现。左侧环境用于前 1000 步，右侧环境用于其余时间。Dyna-Q+ 是带有鼓励探索的探索奖励的 Dyna-Q。}
\end{figure}

当环境改变得比以前更好，而原先正确的策略并不会暴露这种改进时，会出现更大的困难。
在这些情况下，建模错误可能很长时间都检测不到，甚至永远检测不到。

\noindent\textbf{例 8.3：捷径迷宫}
这类环境变化造成的问题，由图 8.5 所示的迷宫例子说明。
最初，最优路径是绕到障碍左侧（左上）。
然而，3000 步之后，在不干扰较长路径的情况下，右侧打开了一条更短的路径（右上）。
图中曲线显示，普通 Dyna-Q 智能体从未切换到这条捷径。
事实上，它从来没有意识到捷径存在。
它的模型认为没有捷径，因此它规划得越多，就越不可能向右走并发现捷径。
即使使用 $\varepsilon$-贪心策略，智能体采取足够多探索动作来发现捷径的可能性也很小。

\begin{figure}[htbp]
\centering
\includegraphics[width=0.72\textwidth]{tex/强化学习/figures/math-9/figure-8-5.png}
\caption{Dyna 智能体在捷径任务上的平均表现。左侧环境用于前 3000 步，右侧环境用于其余时间。}
\end{figure}

这里的一般问题，是探索与利用冲突的另一种形式。
在规划语境中，探索意味着尝试那些能够改进模型的动作，而利用意味着在给定当前模型的情况下以最优方式行动。
我们希望智能体探索，以便发现环境中的变化，但又不希望探索太多，以免表现严重下降。
与前面讨论的探索/利用冲突一样，可能不存在既完美又实用的解法，但简单启发式常常是有效的。

确实解决捷径迷宫的 Dyna-Q+ 智能体，就使用了这样一种启发式。
这个智能体会为每个状态--动作对记录：自从该对上一次在与环境的真实交互中被尝试以来，已经过去了多少时间步。
过去的时间越长，（我们可以假定）该对的动态已经改变、其模型已经不正确的可能性就越大。
为了鼓励测试长期未尝试的动作，在涉及这些动作的模拟经验上，会给出一个特殊的“奖励奖励”。
具体地说，如果某个转移的模型奖励为 $r$，且该转移已有 $\tau$ 个时间步没有被尝试，那么规划更新就好像该转移产生了
\[
r+\kappa\sqrt{\tau}
\]
的奖励，其中 $\kappa$ 是某个较小常数。
这会鼓励智能体持续测试所有可达状态转移，甚至为了进行这些测试而找到很长的动作序列。\footnote{Dyna-Q+ 智能体还做了另外两处修改。第一，在表格型 Dyna-Q 算法的规划步骤 (f) 中，从某个状态以前从未尝试过的动作也允许被考虑。第二，对这些动作的初始模型是：它们会以奖励 0 回到同一状态。}
当然，所有这些测试都有代价；但在许多情况下，如捷径迷宫中，这种计算上的好奇心很值得额外探索。

\noindent\textbf{练习 8.2}
为什么带探索奖励的 Dyna 智能体 Dyna-Q+，在阻塞实验和捷径实验的第一阶段以及第二阶段都表现得更好？

\noindent\textbf{练习 8.3}
仔细查看图 8.5 会发现，在实验的第一部分，Dyna-Q+ 与 Dyna-Q 之间的差距略有缩小。
这是什么原因？

\noindent\textbf{练习 8.4（编程）}
上面描述的探索奖励实际上改变了状态和动作的估计价值。
这是必要的吗？
假设奖励 $\kappa\sqrt{\tau}$ 不用于更新，而只用于动作选择。
也就是说，假设总是选择使
\[
Q(S_t,a)+\kappa\sqrt{\tau(S_t,a)}
\]
最大的动作。
请进行一个网格世界实验，测试并说明这种替代方法的优点和弱点。

\noindent\textbf{练习 8.5}
第 164 页所示的表格型 Dyna-Q 算法可以怎样修改，以处理随机环境？
这种修改在本节考虑的变化环境中可能会怎样表现不佳？
算法又可以怎样修改，以同时处理随机环境和变化环境？

\subsection*{8.4\quad 优先扫掠}

在前几节介绍的 Dyna 智能体中，模拟转移从所有以前经历过的状态--动作对中均匀随机选择而开始。
但均匀选择通常不是最好的；如果把模拟转移和更新集中到特定状态--动作对上，规划可以高效得多。
例如，考虑
第一个迷宫任务（图 8.3）第二幕中发生的事情。
在第二幕开始时，只有直接通向目标的状态--动作对具有正价值；所有其他状态--动作对的价值仍是零。
这意味着，沿着几乎所有转移进行更新都是无意义的，因为这些转移只是把智能体从一个零价值状态带到另一个零价值状态，因此更新不会产生效果。
只有沿着进入目标前一个状态的转移进行更新，或者从这个状态出发的转移进行更新，才会改变任何价值。
如果均匀地产生模拟转移，那么在偶然遇到这些有用更新之前，会做许多浪费的更新。
随着规划推进，有用更新的区域会扩大，但如果规划集中在最能发挥作用的地方，它仍会高效得多。
在我们真正关心的更大问题中，状态数如此之多，以至于无焦点的搜索会极其低效。

这个例子表明，从目标状态\textbf{向后}工作，可能是聚焦搜索的一种有用方式。
当然，我们并不真的希望使用任何专门依赖“目标状态”概念的方法。
我们希望方法适用于一般奖励函数。
目标状态只是一个方便直觉的特例。
一般来说，我们不仅要从目标状态向后工作，还要从任何价值已经改变的状态向后工作。
假设给定模型时，价值最初是正确的，就像迷宫例子在发现目标之前那样。
现在假设智能体发现环境中有变化，并把某个状态的估计价值向上或向下改了。
通常，这意味着许多其他状态的价值也应该改变，但唯一有用的一步更新，是那些直接通向这个价值已改变状态的动作。
如果这些动作的价值被更新，那么它们的前驱状态的价值也可能随之改变。
如果确实如此，那么通向这些前驱状态的动作也需要被更新，然后\textbf{它们}的前驱状态也可能已经改变。
以这种方式，可以从任意价值已改变的状态向后工作，要么执行有用更新，要么终止传播。
这个一般思想可以称为规划计算的\textbf{向后聚焦}。

当有用更新的前沿向后传播时，它常常迅速增长，产生许多可以被有用更新的状态--动作对。
但这些更新并不都同样有用。
某些状态的价值可能改变很多，而另一些改变很少。
那些改变很大的状态，其前驱对也更可能发生较大改变。
在随机环境中，估计转移概率的变化也会造成改变幅度和更新紧迫性的差异。
很自然地，可以按照某种紧迫性度量来给更新排序，并按优先级执行它们。
这就是\textbf{优先扫掠}背后的思想。
维护一个队列，队列中包含每个如果被更新就会使估计价值非平凡改变的状态--动作对，并按改变大小给出优先级。
当队列顶部的对被更新后，就计算这个更新对它所有前驱对的影响。
如果影响大于某个小阈值，那么把该前驱对以新的优先级插入队列；如果该对已经在队列中，则插入后只保留较高优先级的条目。
这样，变化的影响就能高效地向后传播，直到趋于静止。
确定性环境下的完整算法如下框所示。

\begin{center}
\fbox{\begin{minipage}{0.90\textwidth}
\noindent\textbf{确定性环境的优先扫掠}

\vspace{0.4em}
\noindent 对所有 $s,a$ 初始化 $Q(s,a)$、$Model(s,a)$，并令 $PQueue$ 为空。

\noindent 永远循环：

\quad (a) $S\leftarrow$ 当前（非终止）状态；

\quad (b) $A\leftarrow policy(S,Q)$；

\quad (c) 采取动作 $A$；观察得到的奖励 $R$ 和状态 $S'$；

\quad (d) $Model(S,A)\leftarrow R,S'$；

\quad (e) $P\leftarrow |R+\gamma\max_a Q(S',a)-Q(S,A)|$；

\quad (f) 如果 $P>\theta$，则把 $S,A$ 以优先级 $P$ 插入 $PQueue$；

\quad (g) 当 $PQueue$ 非空时，重复 $n$ 次：

\qquad $S,A\leftarrow first(PQueue)$；

\qquad $R,S'\leftarrow Model(S,A)$；

\qquad $Q(S,A)\leftarrow Q(S,A)+\alpha\left[R+\gamma\max_a Q(S',a)-Q(S,A)\right]$；

\qquad 对所有被预测会通向 $S$ 的 $\bar S,\bar A$ 循环：

\qquad\quad $\bar R\leftarrow$ 对 $\bar S,\bar A,S$ 的预测奖励；

\qquad\quad $P\leftarrow |\bar R+\gamma\max_a Q(S,a)-Q(\bar S,\bar A)|$；

\qquad\quad 若 $P>\theta$，则把 $\bar S,\bar A$ 以优先级 $P$ 插入 $PQueue$。
\end{minipage}}
\end{center}

\noindent\textbf{例 8.4：迷宫上的优先扫掠}
已经发现，优先扫掠会显著提高在迷宫任务中找到最优解的速度，常常快 5 到 10 倍。
右图给出了一个典型例子。
这些数据来自一系列迷宫任务，它们与图 8.2 所示任务结构完全相同，只是网格分辨率不同。
相对于未加优先级的 Dyna-Q，优先扫掠保持了决定性优势。
两个系统每次环境交互最多都进行 $n=5$ 次更新。
该例改编自 Peng and Williams（1993）。

\begin{figure}[htbp]
\centering
\includegraphics[width=0.52\textwidth]{tex/强化学习/figures/math-9/figure-8-ps-mazes.png}
\caption{迷宫任务中，优先扫掠与 Dyna-Q 达到最优解之前所需更新次数的比较。}
\end{figure}

把优先扫掠扩展到随机环境是直接的。
模型通过记录每个状态--动作对被经历的次数，以及随后出现过哪些下一个状态来维护。
于是自然的做法是，对每个对不再用我们到目前为止一直使用的样本更新，而是用期望更新，同时考虑所有可能的下一个状态及其发生概率。

优先扫掠只是分配计算以提高规划效率的一种方式，而且很可能不是最好的方式。
优先扫掠的一个限制是它使用期望更新；在随机环境中，这可能会在低概率转移上浪费大量计算。
如下一节所示，尽管抽样会引入方差，样本更新在许多情况下仍能用更少计算更接近真实价值函数。
样本更新可能获胜，是因为它把总体备份计算分解成更小的部分，即对应于单个转移的部分，这样就能更狭窄地集中到影响最大的部分上。
这个思想在 van Seijen and Sutton（2013）引入的“小备份”中被推进到或许是其逻辑极限。
这些更新沿着单个转移进行，像样本更新一样，但它们基于该转移的概率而不进行抽样，就像期望更新一样。
通过选择小备份执行的顺序，可以把规划效率大幅提高到优先扫掠所能达到的程度之上。

\begin{figure}[htbp]
\centering
\includegraphics[width=0.86\textwidth]{tex/强化学习/figures/math-9/example-8-5-rod.png}
\caption{例 8.5 中的杆操纵任务：在带障碍的有限矩形工作空间中，将杆移动到目标位置。}
\end{figure}

本章已经提示，所有状态空间规划都可以看作价值更新序列；它们只是在更新类型（期望或样本、大或小）以及更新顺序上有所不同。
本节强调了向后聚焦，但这只是其中一种策略。
例如，另一种策略是按照状态从当前策略下经常访问的状态到达的难易程度来聚焦，这可以称为\textbf{向前聚焦}。
Peng and Williams（1993）以及 Barto, Bradtke and Singh（1995）探索了向前聚焦的一些版本，后面几节介绍的方法则把它推到一种极端形式。

\subsection*{8.5\quad 期望更新与样本更新}

前几节的例子展示了结合学习方法和规划方法的可能范围。
在本章余下部分，我们分析其中若干组成思想，首先考察期望更新和样本更新的相对优势。

本书很大一部分都在讨论不同类型的价值函数更新，我们已经考虑过许多变体。
暂时聚焦于一步更新，它们主要沿三个二元维度变化。
前两个维度是：更新的是状态价值还是动作价值；估计的是最优策略的价值，还是某个任意给定策略的价值。
这两个维度产生四类更新，用于近似四个价值函数 $q_*,v_*,q_\pi,v_\pi$。
另一个二元维度是：更新是\textbf{期望更新}，考虑所有可能发生的事件；还是\textbf{样本更新}，只考虑可能发生事件的一个单一样本。
这三个二元维度产生八种情形，其中七种对应于具体算法，如右图所示；第八种似乎不对应任何有用更新。
任何这些一步更新都可以用于规划方法。
前面讨论的 Dyna-Q 智能体使用 $q_*$ 样本更新，但它们同样可以使用 $q_*$ 期望更新，或者使用 $q_\pi$ 的期望更新或样本更新。
Dyna-AC 系统把 $v_\pi$ 样本更新与一种学习中的策略结构结合起来使用（如第 13 章）。
对于随机问题，优先扫掠总是使用某种期望更新。

\begin{figure}[htbp]
\centering
\includegraphics[width=0.50\textwidth]{tex/强化学习/figures/math-9/figure-8-6.png}
\caption{本书中考虑的所有一步更新的备份图。}
\end{figure}

第 6 章介绍一步样本更新时，我们把它们作为期望更新的替代物来呈现。
在没有分布模型时，期望更新是不可能的，但样本更新可以使用来自环境或样本模型的样本转移来完成。
这种观点隐含着：如果期望更新可行，那么它们优于样本更新。
但事实真是如此吗？
期望更新当然会给出更好的估计，因为它们不受抽样误差破坏；但它们也需要更多计算，而计算常常是规划中的限制资源。
为了恰当地评价期望更新和样本更新在规划中的相对优劣，我们必须控制它们不同的计算需求。

具体起见，考虑用于近似 $q_*$ 的期望更新和样本更新；并考虑离散状态与动作、近似价值函数 $Q$ 的表格表示，以及估计动态形式为 $\hat p(s',r|s,a)$ 的模型这一特殊情形。
对状态--动作对 $s,a$ 的期望更新是
\begin{equation}
Q(s,a)\leftarrow \sum_{s',r}\hat p(s',r|s,a)\left[r+\gamma\max_{a'}Q(s',a')\right].
\tag{8.1}
\end{equation}
给定来自模型的一个样本下一个状态与奖励 $S'$ 和 $R$，相应的样本更新是类似 Q 学习的更新：
\begin{equation}
Q(s,a)\leftarrow Q(s,a)+\alpha\left[R+\gamma\max_{a'}Q(S',a')-Q(s,a)\right],
\tag{8.2}
\end{equation}
其中 $\alpha$ 是通常的正步长参数。

这些期望更新和样本更新之间差别的显著程度，取决于环境的随机性；更具体地说，取决于在给定状态和动作时，可能有多少下一个状态以各种概率发生。
如果只有一个下一个状态可能发生，那么上面的期望更新和样本更新是相同的（取 $\alpha=1$）。
如果可能有很多下一个状态，那么差别可能很大。
支持期望更新的一点是，它是一种精确计算，产生的新 $Q(s,a)$ 的正确性只受后继状态 $Q(s',a')$ 正确性的限制。
样本更新还会额外受到抽样误差影响。
另一方面，样本更新计算上更便宜，因为它只考虑一个下一个状态，而不是所有可能的下一个状态。
在实践中，更新操作所需的计算通常主要由需要求值 $Q$ 的状态--动作对数量决定。
对某个起始对 $s,a$，令 $b$ 为\textbf{分支因子}，也就是满足 $\hat p(s'|s,a)>0$ 的可能下一个状态 $s'$ 的个数。
那么，对这个对进行一次期望更新，大约需要一次样本更新的 $b$ 倍计算。

如果有足够时间完成一次期望更新，那么由于没有抽样误差，得到的估计通常优于 $b$ 次样本更新。
但如果没有足够时间完成一次期望更新，那么样本更新总是更可取，因为它们至少能用少于 $b$ 次更新在价值估计上取得一些改进。
在有许多状态--动作对的大问题中，我们处于后一种情况的概率很大。
在这种情况下，若要对所有这些对进行期望更新，会花很长时间。
在那之前，对许多状态--动作对进行少量样本更新，可能远好于只对少数对进行期望更新。
给定一个单位的计算努力，它究竟更应该用于少量期望更新，还是用于 $b$ 倍数量的样本更新？

图 8.7 显示了一个分析结果，为这个问题提供了答案的线索。
它显示了在多种分支因子 $b$ 下，期望更新和样本更新的估计误差随计算时间变化的情况。
考虑的情形是：$b$ 个后继状态等可能发生，且初始估计误差为 1。
假定下一个状态的价值是正确的，所以期望更新一完成就把误差降为零。
在这种情况下，样本更新会按照
\[
\sqrt{\frac{b-1}{bt}}
\]
减小误差，其中 $t$ 是已经执行的样本更新次数（假设使用样本平均，即 $\alpha=1/t$）。
关键观察是，对于中等偏大的 $b$，误差在只进行了 $b$ 的很小一部分更新后就急剧下降。
在这些情形下，许多状态--动作对的价值都可以在一次状态--动作对完成期望更新所需的同样时间内，被显著改进到期望更新效果的几个百分点之内。

\begin{figure}[htbp]
\centering
\includegraphics[width=0.82\textwidth]{tex/强化学习/figures/math-9/figure-8-7.png}
\caption{期望更新与样本更新效率的比较。}
\end{figure}

图 8.7 中显示的样本更新优势，可能低估了真实效果。
在真实问题中，后继状态的价值本身也会是正在被更新的估计。
由于样本更新会使估计更快变得更准确，所以样本更新还会有第二个优势：从后继状态备份而来的价值也会更准确。
这些结果表明，对于随机分支因子很大、且状态太多而无法精确求解的问题，样本更新很可能优于期望更新。

\noindent\textbf{练习 8.6}
上面的分析假定 $b$ 个可能的下一个状态都等可能发生。
如果分布高度偏斜，使得 $b$ 个状态中有些比大多数状态更可能发生，会怎样？
这会加强还是削弱样本更新相对于期望更新的理由？
请支持你的回答。

\subsection*{8.6\quad 轨迹采样}

本节比较两种分配更新的方法。
动态规划中的经典方法是扫过整个状态（或状态--动作）空间，每一扫对每个状态（或状态--动作对）更新一次。
这在
大任务中有问题，因为可能没有时间完成哪怕一遍扫掠。
在许多任务中，绝大多数状态都是无关的，因为它们只会在很差的策略下被访问，或者被访问的概率非常低。
穷尽扫掠隐含地把相同时间分给状态空间的所有部分，而不是集中到真正需要的地方。
正如第 4 章所讨论的，穷尽扫掠及其对所有状态的平等处理，并不是动态规划的必要性质。
原则上可以按任意方式分配备份（为了保证收敛，所有状态或状态--动作对在极限中必须被访问无穷多次；不过第 8.7 节会讨论一个例外），但实践中仍常使用穷尽扫掠。

第二种方法是按照某个分布从状态或状态--动作空间中抽样。
可以像 Dyna-Q 智能体那样均匀抽样，但这会遭遇与穷尽扫掠相同的一些问题。
更有吸引力的是按照\textbf{同策略分布}来分配更新，也就是按照遵循当前策略时观察到的分布。
这个分布的一个优点是很容易生成；只需遵循当前策略与模型交互即可。
在分幕任务中，从起始状态（或按起始状态分布）开始并模拟直到终止状态。
在持续性任务中，从任何地方开始并不断模拟。
无论哪种情况，样本状态转移和奖励由模型给出，样本动作由当前策略给出。
换句话说，我们模拟显式的单条轨迹，并对沿途遇到的状态或状态--动作对进行更新。
我们把这种生成经验并更新的方式称为\textbf{轨迹采样}。

除了轨迹采样，很难想象还有什么高效方法能按照同策略分布分配更新。
如果有同策略分布的显式表示，那么可以扫过所有状态，并按同策略分布为每个状态的更新加权，但这又回到了穷尽扫掠的全部计算代价。
也许可以从这个分布中抽样并更新单个状态--动作对，但即便能高效做到，相比模拟轨迹又有什么好处呢？
甚至知道同策略分布的显式形式也不太可能。
每当策略改变时，这个分布就会改变，而计算该分布需要与完整策略评价相当的计算量。
考虑这些其他可能性，会使轨迹采样显得既高效又优雅。

同策略更新分布好吗？
直觉上，它似乎是一个好选择，至少好于均匀分布。
例如，如果你正在学习下棋，就会研究真实对局中可能出现的局面，而不是随机摆放棋子的局面。
后者也许是合法状态，但准确评价它们是一种不同于评价真实对局局面的技能。
我们还会在第二部分看到，在使用函数近似时，同策略分布有显著优势。
无论是否使用函数近似，人们都可以期待同策略聚焦会显著提高规划速度。

聚焦于同策略分布可能有利，因为它会使空间中巨大而无趣的部分被忽略；也可能有害，因为它会使同一些旧区域被一再更新。
我们进行了一个小实验来经验评估其影响。
为了隔离更新分布的影响，我们完全使用式 (8.1) 定义的一步期望表格更新。
在\textbf{均匀}情形中，我们循环遍历所有状态--动作对，并原地更新每一个。
在\textbf{同策略}情形中，我们模拟幕，所有幕都从同一状态开始，并更新当前 $\varepsilon$-贪心策略（$\varepsilon=0.1$）下出现的每个状态--动作对。
这些任务是不折扣的分幕任务，随机生成如下。
对 $|\mathcal{S}|$ 个状态中的每一个，都有两个可能动作；每个动作会导致 $b$ 个下一个状态之一，这些状态等可能发生，并且每个状态--动作对对应的 $b$ 个状态都是不同的随机选择。
分支因子 $b$ 对所有状态--动作对都相同。
此外，在所有转移上，都有 0.1 的概率转移到终止状态，从而结束该幕。
每个转移上的期望奖励从均值 0、方差 1 的高斯分布中选出。
在规划过程的任一点，都可以停止并穷尽计算 $v_{\tilde\pi}(s_0)$：在当前动作价值函数 $Q$ 所诱导的贪心策略 $\tilde\pi$ 下，起始状态的真实价值，以此作为智能体在一个新幕中贪心行动时表现如何的指标（同时假设模型是正确的）。

图 8.8 的上半部分显示了 1000 个状态、分支因子为 1、3 和 10 的 200 个样本任务上的平均结果。
图中把找到的策略质量画成已完成期望更新次数的函数。
在所有情形中，按照同策略分布抽样在初期都带来更快规划，而长期来看会拖慢规划。
分支因子越小，这种效果越强，初始更快规划的时期也越长。
在其他实验中，我们发现这些效果也会随着状态数增加而变强。
例如，图的下半部分显示了 10000 个状态且分支因子为 1 的任务结果。
在这种情况下，同策略聚焦的优势很大，而且持续很久。

\begin{figure}[htbp]
\centering
\includegraphics[width=0.58\textwidth]{tex/强化学习/figures/math-9/figure-8-8.png}
\caption{均匀分布于状态空间的更新，与集中在模拟同策略轨迹上的更新之间的相对效率比较；每条轨迹都从同一状态开始。结果来自两种规模和多种分支因子 $b$ 的随机生成任务。}
\end{figure}

所有这些结果都合乎情理。
短期来看，按照同策略分布抽样有帮助，因为它聚焦于起始状态的近后代状态。
如果状态很多且分支因子很小，这种效果会很大并持续很久。
长期来看，聚焦于同策略分布可能有害，因为常出现的状态已经具有正确价值。
继续抽样它们是无用的，而抽样其他状态也许实际上能完成一些有用工作。
这大概就是为什么穷尽的、无聚焦的方法在长期中表现更好，至少在小问题中如此。
这些结果并非结论性的，因为它们只针对一种特定随机方式生成的问题；但它们确实表明，对大问题而言，按照同策略分布抽样可以带来巨大优势，特别是当同策略分布只访问状态--动作空间中一个小子集时。

\noindent\textbf{练习 8.7}
图 8.8 中有些曲线在早期部分似乎呈扇贝形，特别是 $b=1$ 且均匀分布的上图。
你认为为什么会这样？
图中数据的哪些方面支持你的假设？

\noindent\textbf{练习 8.8（编程）}
复现图 8.8 下半部分所示实验，然后尝试同一个实验但令 $b=3$。
讨论你的结果含义。

\subsection*{8.7\quad 实时动态规划}

\textbf{实时动态规划}（real-time dynamic programming，RTDP）是动态规划（DP）价值迭代算法的同策略轨迹采样版本。
因为它与传统的基于扫掠的策略迭代关系密切，所以 RTDP 特别清楚地说明了同策略轨迹采样能提供的一些优点。
RTDP 通过式 (4.10) 定义的期望表格型价值迭代更新，来更新实际或模拟轨迹中访问到的状态价值。
它基本上就是产生图 8.8 中同策略结果的算法。

RTDP 与传统 DP 之间的紧密联系，使我们能够通过改造已有理论来得到若干理论结果。
RTDP 是第 4.5 节所述的\textbf{异步 DP} 算法的一个例子。
异步 DP 算法不是按状态集的系统扫掠组织起来的；它们以任意顺序更新状态价值，并使用当时可用的其他状态价值。
在 RTDP 中，更新顺序由真实或模拟轨迹中状态被访问的顺序决定。

如果轨迹只能从指定的起始状态集合开始，并且你关心的是给定策略下的预测问题，那么同策略轨迹采样允许算法完全跳过那些无法由给定策略从任何起始状态到达的状态；这些状态与预测问题\textbf{无关}。
对于控制问题，目标是寻找最优策略而不是评价给定策略，那么很可能也有一些状态无法由任何最优策略从任何起始状态到达，而且没有必要为这些无关状态指定最优动作。
需要的是一个\textbf{最优偏策略}，意思是它在相关状态上是最优的，但在无关状态上可以指定任意动作，甚至可以未定义。

\begin{figure}[htbp]
\centering
\includegraphics[width=0.55\textwidth]{tex/强化学习/figures/math-9/rtdp-relevant-states.png}
\caption{RTDP 中相关状态与无关状态的示意。相关状态是从某个起始状态出发、在某个最优策略下可达的状态。}
\end{figure}

但是，用同策略轨迹采样控制方法寻找这样的最优偏策略，例如 Sarsa（第 6.4 节），通常需要访问所有状态--动作对，甚至那些最终会变得无关的对，也要访问无穷多次。
例如，可以使用探索起始（第 5.3 节）做到这一点。
RTDP 也是如此：对于带探索起始的分幕任务，RTDP 是一种异步价值迭代算法，会收敛到折扣有限 MDP 的最优策略（在某些条件下，也会在不折扣情形中收敛）。
与预测问题不同，如果收敛到最优策略很重要，那么通常不可能停止更新任何状态或状态--动作对。

RTDP 最有趣的结果是：对满足合理条件的某些问题类型，RTDP 被保证能找到一个在相关状态上最优的策略，而不必无穷多次访问每个状态，甚至不必访问某些状态。
事实上，在某些问题中，只需要访问很小一部分状态。
对于状态集非常大的问题，这是一个巨大优势，因为甚至单次扫掠都可能不可行。

这个结果适用的任务，是具有吸收目标状态且这些目标状态产生零奖励的 MDP 上的不折扣分幕任务，如第 3.4 节所述。
在真实或模拟轨迹的每一步，RTDP 选择一个贪心动作（随机打破平局），并对当前状态应用期望价值迭代更新操作。
它也可以在每一步更新任意一批其他状态的价值；例如，可以更新从当前状态出发的有限视界前瞻搜索中访问到的状态价值。

对这些问题，当每一幕从起始状态集合中随机选择的状态开始并在目标状态结束时，若满足以下条件，RTDP 以概率 1 收敛到一个对所有相关状态都最优的策略：
1）每个目标状态的初始价值为零；
2）至少存在一个策略，能保证从任一起始状态以概率 1 到达某个目标状态；
3）所有从非目标状态出发的转移奖励都严格为负；
4）所有初始价值都等于或大于其最优价值（只需把所有状态的初始价值设为零即可满足）。
这个结果由 Barto, Bradtke, and Singh（1995）证明；他们把异步 DP 的结果与 Korf（1990）提出的启发式搜索算法 learning real-time A* 的结果结合起来。

具有这些性质的任务，是\textbf{随机最优路径问题}的例子。
这类问题通常以成本最小化表述，而不是像这里这样以奖励最大化表述。
在我们的版本中最大化负回报，等价于最小化从起始状态到目标状态的路径成本。
这类任务的例子包括最短时间控制任务，其中到达目标前每个时间步产生奖励 $-1$；也包括第 3.5 节的高尔夫例子，其目标是用最少杆数把球打入洞中。

\noindent\textbf{例 8.6：赛车场上的 RTDP}
练习 5.12（第 111 页）的赛车场问题是一个随机最优路径问题。
把 RTDP 与传统 DP 价值迭代算法在一个赛车场例题上比较，可以说明同策略轨迹采样的一些优点。

回忆该练习：智能体必须学会如何驾驶汽车绕过类似图 5.5 所示的弯道，并在保持在赛道上的同时尽快越过终点线。
起始状态是起跑线上所有速度为零的状态；目标状态是所有能够通过从赛道内部越过终点线而在一个时间步内到达的状态。
与练习 5.12 不同，这里对汽车速度没有限制，所以状态集潜在地是无限的。
不过，从起始状态集合出发、经由任意策略能够到达的状态集合是有限的，可以把它看作问题的状态集。
每一幕从随机选择的起始状态开始，并在汽车越过终点线时结束。
每一步直到汽车越过终点线之前，奖励都是 $-1$。
如果汽车碰到赛道边界，它会被移回一个随机起始状态，幕继续进行。

一个类似于图 5.5 左侧小赛车场的赛道，从起始状态出发经任意策略可达的状态有 9115 个，其中只有 599 个是相关的，也就是可由某个最优策略从某个起始状态到达。
相关状态数通过统计执行最优动作 $10^7$ 幕期间访问到的状态估计得到。

下表比较了用传统 DP 和 RTDP 求解该任务。
这些结果是在 25 次运行上取平均得到的，每次运行用不同随机数种子开始。
这里的传统 DP 是价值迭代，使用对状态集的穷尽扫掠，每次原地更新一个状态的价值，这意味着每个状态的更新都使用其他状态的最新价值。
这就是价值迭代的 Gauss--Seidel 版本；在这个问题上，它被发现约为 Jacobi 版本的两倍快（见第 4.8 节）。
没有特别注意更新顺序；其他顺序可能会产生更快收敛。
两种方法每次运行的初始价值都全为零。
当某次扫掠中状态价值的最大变化小于 $10^{-4}$ 时，判定 DP 已收敛；当 20 幕上越过终点线的平均时间看起来稳定在渐近步数时，判定 RTDP 已收敛。
这个版本的 RTDP 在每一步只更新当前状态的价值。

\begin{center}
\begin{tabular}{lcc}
\toprule
 & DP & RTDP \\
\midrule
收敛所需平均计算量 & 28 次扫掠 & 4000 幕 \\
收敛所需平均更新次数 & 252,784 & 127,600 \\
每幕平均更新次数 & -- & 31.9 \\
更新不超过 100 次的状态百分比 & -- & 98.45 \\
更新不超过 10 次的状态百分比 & -- & 80.51 \\
更新 0 次的状态百分比 & -- & 3.18 \\
\bottomrule
\end{tabular}
\end{center}

两种方法产生的策略平均都需要 14 到 15 步越过终点线，但 RTDP 所需更新次数大约只有 DP 的一半。
这是 RTDP 的同策略轨迹采样造成的。
在 DP 的每次扫掠中，每个状态的价值都会被更新；而 RTDP 把更新集中在更少状态上。
一次平均运行中，RTDP 对 98.45\% 的状态更新不超过 100 次，对 80.51\% 的状态更新不超过 10 次；平均约有 290 个状态的价值根本没有被更新。

RTDP 的另一个优点是，随着价值函数接近最优价值函数 $v_*$，智能体用来生成轨迹的策略也接近最优策略，因为它总是相对于当前价值函数贪心。
这与传统价值迭代中的情况形成对比。
实践中，价值迭代会在一次扫掠中价值函数只发生很小变化时终止；上表中的结果也是这样终止的。
此时价值函数近似 $v_*$，贪心策略接近最优。
然而，完全可能在价值迭代终止很久之前，相对于最新价值函数的贪心策略已经是最优或近似最优的（回忆第 4 章：最优策略可以相对于许多不同价值函数贪心，而不只是 $v_*$）。
在价值迭代收敛之前检查最优策略是否出现，并不是传统 DP 算法的一部分，而且需要相当多的额外计算。

在赛车场例子中，通过每次 DP 扫掠后运行许多测试幕，并按照该次扫掠结果贪心选择动作，可以估计 DP 计算中最早何时近似最优价值函数已经足够好，以至于相应的贪心策略近似最优。
对这个赛车场而言，一个接近最优的策略在 15 次价值迭代扫掠之后出现，也就是在 136,725 次价值迭代更新之后出现。
这显著少于 DP 收敛到 $v_*$ 所需的 252,784 次更新，但仍多于 RTDP 所需的 127,600 次更新。

虽然这些模拟当然不是 RTDP 与传统基于扫掠价值迭代之间的决定性比较，但它们说明了同策略轨迹采样的一些优点。
传统价值迭代继续更新所有状态的价值，而 RTDP 强烈聚焦于与问题目标相关的状态子集。
随着学习继续，这种聚焦越来越窄。
由于 RTDP 的收敛定理适用于这些模拟，我们知道 RTDP 最终只会聚焦于相关状态，也就是构成最优路径的状态。
RTDP 用大约基于扫掠的价值迭代 50\% 的计算量，实现了近似最优控制。

\subsection*{8.8\quad 决策时规划}

规划至少可以用两种方式。
本章到目前为止考虑的一种方式，以动态规划和 Dyna 为代表：使用规划，根据从模型（样本模型或分布模型）得到的模拟经验，逐渐改进策略或价值函数。
随后，动作选择就是比较当前状态的动作价值；在目前考虑的表格情形中，这些动作价值来自表，而在第二部分考虑的近似方法中，则来自对某个数学表达式的求值。
早在为任何当前状态 $S_t$ 选择动作之前，规划已经在改进表项或
数学表达式，它们随后会用于包括 $S_t$ 在内的许多状态上的动作选择。
以这种方式使用时，规划并不聚焦于当前状态。
我们称这种规划为\textbf{后台规划}。

使用规划的另一种方式，是在每次遇到新状态 $S_t$ 之后开始并完成一次规划计算；这次计算的输出是单个动作 $A_t$ 的选择。
下一步中，规划又从 $S_{t+1}$ 重新开始，以产生 $A_{t+1}$，依此类推。
这种规划最简单、也几乎退化的例子，是当只有状态价值可用时，通过比较每个动作下模型预测下一个状态的价值来选择动作（或像第 1 章井字棋例子那样，比较后状态的价值）。
更一般地，以这种方式使用的规划可以看得比一步前瞻深得多，并评价会通向许多不同预测状态与奖励轨迹的动作选择。
与规划的第一种用法不同，这里规划聚焦于某个特定状态。
我们称这种规划为\textbf{决策时规划}。

这两种关于规划的思考方式，即使用模拟经验逐渐改进策略或价值函数，或者使用模拟经验为当前状态选择动作，可以以自然且有趣的方式混合起来。
不过，它们往往被分开研究，而这也是最初理解它们的好方式。
现在我们更仔细地考察决策时规划。

即使规划只在决策时进行，我们仍然可以像第 8.1 节那样，把它看成从模拟经验出发，经由更新和价值，最终到达策略。
只是现在价值和策略都特定于当前状态以及那里可用的动作选择，以至于规划过程产生的价值和策略通常在被用来选择当前动作之后就被丢弃。
在许多应用中，这并不是很大的损失，因为状态非常多，我们不太可能很快回到同一状态。
一般来说，人们可能希望二者混合：把规划聚焦于当前状态，同时也保存规划结果，以便将来若回到同一状态时能走得更远。
决策时规划在不要求快速响应的应用中最有用。
例如，在国际象棋程序中，每步棋可以允许数秒或数分钟计算，强程序能在这段时间内向前规划几十步。
另一方面，如果低延迟动作选择是优先目标，那么通常更适合在后台进行规划，计算出一个随后可以快速应用于每个新遇到状态的策略。

\subsection*{8.9\quad 启发式搜索}

人工智能中经典的状态空间规划方法，是统称为\textbf{启发式搜索}的决策时规划方法。
在启发式搜索中，对遇到的每个状态，都考虑一棵很大的可能后续情形树。
近似价值函数被应用于叶节点，然后向搜索树根部的当前状态备份。
搜索树内部的备份，与本书一直讨论的带最大化的期望更新（即用于 $v_*$ 和 $q_*$ 的更新）完全相同。
备份在当前状态的状态--动作节点处停止。
一旦这些节点的备份值被计算出来，其中最好的一个就被选为当前动作，然后所有备份值都被丢弃。

在传统启发式搜索中，并不试图通过改变近似价值函数来保存这些备份值。
事实上，价值函数通常由人设计，并不会因搜索而改变。
然而，让价值函数随时间改进是自然的；可以使用启发式搜索期间计算出的备份值，也可以使用本书中介绍的其他任何方法。
从某种意义上说，我们一直在采用这种方法。
我们的贪心、$\varepsilon$-贪心和 UCB（第 2.7 节）动作选择方法并不像启发式搜索，但在较小规模上是类似的。
例如，给定一个模型和一个状态价值函数，要计算贪心动作，我们必须从每个可能动作向前看到每个可能下一个状态，考虑奖励和估计价值，然后选择最佳动作。
正如传统启发式搜索中那样，这个过程计算了可能动作的备份值，但并不试图保存它们。
因此，启发式搜索可以看作把贪心策略思想从一步扩展到多步。

搜索超过一步的目的，是得到更好的动作选择。
如果有一个完美模型和一个不完美的动作价值函数，那么更深搜索通常确实会产生更好的策略。\footnote{对此也有一些有趣例外；例如见 Pearl, 1984。}
当然，如果搜索一直进行到幕结束，那么不完美价值函数的影响就被消除了，用这种方式确定的动作必然是最优的。
如果搜索深度 $k$ 足够大，使得 $\gamma^k$ 很小，那么动作也会相应地接近最优。
另一方面，搜索越深，需要的计算越多，通常导致响应时间更慢。
Tesauro 的大师级双陆棋程序 TD-Gammon（第 16.1 节）提供了一个好例子。
这个系统通过许多自我对弈学习了一个后状态价值函数，并使用一种启发式搜索来走棋。
作为模型，TD-Gammon 使用关于掷骰概率的先验知识，并假定对手总是选择 TD-Gammon 认为对它最好的动作。
Tesauro 发现，启发式搜索越深，TD-Gammon 走出的棋越好，但每步棋耗时也越长。
双陆棋的分支因子很大，但走棋必须在几秒内完成。
实际可行的只是有选择地向前搜索几步；即便如此，搜索仍显著改善了动作选择。

我们不应忽视启发式搜索聚焦更新的最明显方式：聚焦于当前状态。
启发式搜索有效性的很大一部分，来自它的搜索树紧密集中在当前状态之后可能立即出现的状态和动作上。
你也许一生中下国际象棋比跳棋更多，但当你下跳棋时，思考跳棋以及你当前特定跳棋局面、可能的下一步和后继局面是值得的。
无论你怎样选择动作，这些状态和动作都是更新的最高优先级所在，也是你最迫切希望近似价值函数准确的地方。
不仅计算应优先用于迫近事件，有限的记忆资源也应如此。
例如，在国际象棋中，可能局面太多，不可能为每个局面存储不同价值估计；但基于启发式搜索的国际象棋程序可以很容易地为从单个局面向前看时遇到的数百万个局面存储不同估计。
这种把记忆和计算资源强烈聚焦到当前决策上的方式，大概就是启发式搜索如此有效的原因。

更新分布也可以用类似方式改变，以聚焦于当前状态及其可能后继。
作为极限情形，我们可以完全使用启发式搜索的方法构造搜索树，然后如图 8.9 所示，从底向上执行单个一步更新。
如果更新以某种方式排序，并且使用表格表示，那么可以得到与深度优先启发式搜索完全相同的总体更新。
任何状态空间搜索都可以这样看作由大量单个一步更新拼接而成。
因此，更深搜索带来的性能改进，并不是因为使用了多步更新本身。
相反，它来自更新对当前状态下游紧邻状态和动作的聚焦与集中。
通过把大量计算投入到与候选动作特别相关的内容，决策时规划可以产生优于依赖无焦点更新的决策。

\begin{figure}[htbp]
\centering
\includegraphics[width=0.82\textwidth]{tex/强化学习/figures/math-9/figure-8-9.png}
\caption{启发式搜索可以实现为一系列一步更新；图中用蓝色轮廓表示这些更新，它们把价值从叶节点向根部备份。这里显示的顺序对应于选择性的深度优先搜索。}
\end{figure}

\subsection*{8.10\quad Rollout 算法}

Rollout 算法是基于蒙特卡洛控制的决策时规划算法；这些蒙特卡洛控制应用于都从当前环境状态开始的模拟轨迹。
它们通过平均许多模拟轨迹的回报来估计给定策略的动作价值；这些轨迹分别从每个可能动作开始，然后遵循给定策略。
当动作价值估计被认为足够准确时，就执行估计价值最高的动作（或其中一个动作），然后从得到的下一个状态重新进行整个过程。
正如 Tesauro and Galperin（1997）在将 rollout 算法用于双陆棋时解释的那样，“rollout”一词来自这样一种做法：为了估计双陆棋局面的价值，把该局面用随机生成的掷骰序列多次“展开”到游戏结束，而双方走法由某个固定策略决定。

与第 5 章描述的蒙特卡洛控制算法不同，rollout 算法的目标不是估计完整的最优动作价值函数 $q_*$，也不是估计某个给定策略 $\pi$ 的完整动作价值函数 $q_\pi$。
相反，它们只为每个当前状态、以及一个通常称为\textbf{rollout 策略}的给定策略，产生动作价值的蒙特卡洛估计。
作为决策时规划算法，rollout 算法立即使用这些动作价值估计，然后丢弃它们。
这使 rollout 算法实现起来相对简单，因为没有必要为每个状态--动作对抽样结果，也没有必要在状态空间或状态--动作空间上近似一个函数。

那么 rollout 算法到底实现了什么？
第 4.2 节的策略改进定理告诉我们：给定任意两个策略 $\pi$ 和 $\pi'$，如果它们除了在某个状态 $s$ 上 $\pi'(s)=a\ne\pi(s)$ 之外完全相同，并且 $q_\pi(s,a)\ge v_\pi(s)$，那么策略 $\pi'$ 至少和 $\pi$ 一样好。
此外，如果不等式严格成立，那么 $\pi'$ 实际上优于 $\pi$。
这适用于 rollout 算法，其中 $s$ 是当前状态，$\pi$ 是 rollout 策略。
对模拟轨迹的回报取平均，会为每个动作 $a'\in\mathcal{A}(s)$ 产生 $q_\pi(s,a')$ 的估计。
随后，在 $s$ 中选择使这些估计最大的动作的策略，是改进 rollout 策略 $\pi$ 的一个好候选策略。
结果类似于第 4.3 节讨论的动态规划策略迭代算法中的一步策略改进；不过，由于它只改变当前状态的动作，所以更像第 4.5 节所述异步价值迭代的一步。

换句话说，rollout 算法的目标是改进 rollout 策略，而不是寻找最优策略。
经验表明，rollout 算法可以出人意料地有效。
例如，Tesauro and Galperin（1997）惊讶于 rollout 方法给双陆棋水平带来的巨大提升。
在某些应用中，即使 rollout 策略完全随机，rollout 算法也能产生良好表现。
但改进后策略的表现依赖 rollout 策略的性质，也依赖蒙特卡洛价值估计给出的动作排序。
直觉表明，rollout 策略越好、价值估计越准确，rollout 算法产生的策略就越可能更好（但可参见 Gelly and Silver, 2007）。

这里涉及重要权衡，因为更好的 rollout 策略通常意味着需要更多时间模拟足够多轨迹，以获得良好价值估计。
作为决策时规划方法，rollout 算法通常必须满足严格时间限制。
rollout 算法所需计算时间，取决于每次决策必须评价的动作数、获得有用样本回报所需模拟轨迹的时间步数、rollout 策略做决策所需的时间，以及获得良好蒙特卡洛动作价值估计所需的模拟轨迹数量。

在 rollout 方法的任何应用中，平衡这些因素都很重要，但有几种方式可以缓解挑战。
由于蒙特卡洛试验彼此独立，可以在不同处理器上并行运行许多试验。
另一种方法是把模拟轨迹截断，不展开到完整幕，并用存储的评价函数来校正截断回报（这会引入我们关于截断回报以及前面章节更新的所有讨论）。
还可以像 Tesauro and Galperin（1997）建议的那样，监控蒙特卡洛模拟，并剪掉那些不太可能成为最优的候选动作，或者那些价值已经足够接近当前最优动作、以至于选择它们不会造成实质差异的候选动作（不过 Tesauro 和 Galperin 指出，这会使并行实现复杂化）。

我们通常不把 rollout 算法看作\textbf{学习}算法，因为它们不维护价值或策略的长期记忆。
然而，这些算法利用了本书一直强调的强化学习若干特征。
作为蒙特卡洛控制实例，它们通过平均一组样本轨迹的回报来估计动作价值；在这里，这些轨迹是与环境样本模型交互的模拟轨迹。
以这种方式，它们像强化学习算法一样，通过轨迹采样避免动态规划的穷尽扫掠，并通过依赖样本更新而不是期望更新，避免了对分布模型的需要。
最后，rollout 算法通过相对于估计动作价值贪心地行动，利用了策略改进性质。

\subsection*{8.11\quad 蒙特卡洛树搜索}

\textbf{蒙特卡洛树搜索}（Monte Carlo Tree Search，MCTS）是决策时规划的一个较新且极其成功的例子。
从根本上说，MCTS 是上面描述的 rollout 算法，但通过增加一种机制来累积从蒙特卡洛模拟中得到的价值估计，从而把后续模拟逐渐引向奖励更高的轨迹。
MCTS 在很大程度上促成了计算机围棋从 2005 年的弱业余水平提升到 2015 年的大师水平（6 段或更高）。
基本算法已经发展出许多变体，包括我们将在第 16.6 节讨论的一种变体；它对于 AlphaGo 程序在 2016 年战胜一位 18 次世界冠军围棋棋手具有关键作用。
MCTS 已经证明在各种竞争性场景中有效，包括通用游戏（例如 Finnsson and Björnsson, 2008；Genesereth and Thielscher, 2014），但它并不限于游戏；如果环境模型足够简单，能够快速进行多步模拟，它也能有效用于单智能体序贯决策问题。

MCTS 在遇到每个新状态后执行，为该状态选择智能体的动作；在下一个状态中又再次执行，以选择下一动作，依此类推。
与 rollout 算法一样，每次执行都是一个迭代过程：它模拟许多从当前状态开始并运行到终止状态的轨迹（或者直到折扣使进一步奖励对回报的贡献可以忽略）。
MCTS 的核心思想是，通过
扩展那些在较早模拟中得到高评价的轨迹初始部分，逐步聚焦从当前状态开始的多次模拟。
MCTS 不必从一次动作选择到下一次动作选择都保留近似价值函数或策略，虽然在许多实现中，它会保留下一次执行可能有用的一些动作价值。

在大多数情况下，模拟轨迹中的动作由一个简单策略生成；与更简单的 rollout 算法一样，这个策略通常称为 rollout 策略。
当 rollout 策略和模型都不需要太多计算时，可以在很短时间内生成许多模拟轨迹。
与任何表格型蒙特卡洛方法一样，一个状态--动作对的价值被估计为从该对出发的（模拟）回报平均。
蒙特卡洛价值估计只为最可能在几步内到达的状态--动作对的子集维护；这些对构成一棵以当前状态为根的树，如图 8.10 所示。
MCTS 通过添加那些根据模拟轨迹结果看起来有希望的状态节点，增量扩展这棵树。
任何模拟轨迹都会穿过这棵树，然后在某个叶节点离开。
在树外以及叶节点处，动作选择使用 rollout 策略；但在树内状态上，可以做得更好。
对这些状态，我们已经有至少一部分动作的价值估计，因此可以使用一种有信息的策略在它们之间选择。
这种策略称为\textbf{树策略}，它在探索与利用之间取得平衡。
例如，树策略可以使用 $\varepsilon$-贪心或 UCB 选择规则（第 2 章）来选择动作。

\begin{figure}[htbp]
\centering
\includegraphics[width=0.86\textwidth]{tex/强化学习/figures/math-9/figure-8-10.png}
\caption{蒙特卡洛树搜索。当环境转移到一个新状态时，MCTS 会在必须选择动作之前尽可能多地执行迭代，逐步构建一棵以当前状态为根节点的树。每次迭代包括选择、扩展、模拟和备份四个操作；图中用粗箭头示意。}
\end{figure}

更详细地说，基本 MCTS 的每次迭代由图 8.10 所示的四个步骤组成：
\begin{enumerate}
\item \textbf{选择。} 从根节点开始，一个基于树边上动作价值的树策略遍历这棵树，以选择一个叶节点。
\item \textbf{扩展。} 在某些迭代中（取决于应用细节），从被选中的叶节点扩展树：添加一个或多个从该节点经由尚未探索动作到达的子节点。
\item \textbf{模拟。} 从被选中的节点，或从其新加入的某个子节点（若有）开始，用 rollout 策略选择动作，运行一个完整幕的模拟。其结果是一个蒙特卡洛试验：先由树策略选择动作，离开树之后由 rollout 策略选择动作。
\item \textbf{备份。} 模拟幕产生的回报被备份，用来更新或初始化本次 MCTS 迭代中树策略所遍历树边上的动作价值。对 rollout 策略在树外访问到的状态和动作，不保存价值。图 8.10 通过从模拟轨迹的终止状态直接备份到树中 rollout 策略开始处的状态--动作节点，说明了这一点（不过一般来说，模拟轨迹的整个回报会备份到这个状态--动作节点）。
\end{enumerate}

MCTS 从树的根节点开始反复执行这四步，直到时间耗尽或其他计算资源耗尽。
然后，最后根据树中累积的统计量，从根节点（它仍表示环境的当前状态）选择一个动作；例如，可以选择根状态所有可用动作中动作价值最大的动作，或者为了避免选择离群点而选择访问次数最多的动作。
这就是 MCTS 实际选择的动作。
环境转移到新状态后，再次运行 MCTS；有时从只包含表示新状态的单个根节点的树开始，但常常会从上一轮 MCTS 构造的树中，保留该新节点的所有后代作为起始树；所有其他节点以及与它们关联的动作价值都被丢弃。

MCTS 最初是为二人竞争性游戏程序（如围棋）选择走法而提出的。
在游戏中，每个模拟幕都是一整局棋，其中双方都由树策略和 rollout 策略来选择动作。
第 16.6 节描述了 AlphaGo 程序中使用的 MCTS 扩展，它把 MCTS 的蒙特卡洛评价与通过自我对弈强化学习由深度人工神经网络学到的动作价值结合起来。

把 MCTS 与本书描述的强化学习原则联系起来，有助于理解它为何取得如此惊人的结果。
从根本上说，MCTS 是一种基于蒙特卡洛控制的决策时规划算法；这些蒙特卡洛控制应用于从根状态开始的模拟。
也就是说，它是上一节所述的一类 rollout 算法。
因此，它受益于在线的、增量的、基于样本的价值估计和策略改进。
除此之外，它还保存附着在树边上的动作价值估计，并用强化学习的样本更新来更新它们。
这会把蒙特卡洛试验聚焦到那些初始片段与先前模拟出的高回报轨迹相同的轨迹上。
进一步说，通过增量扩展树，MCTS 实际上增长出一张查找表，用来存储一个偏动作价值函数；内存被分配给高收益样本轨迹初始片段中访问到的状态--动作对的估计价值。
因此，MCTS 避免了全局近似动作价值函数的问题，同时保留了利用过去经验引导探索的好处。

MCTS 的决策时规划取得的显著成功，已经深刻影响人工智能；许多研究者正在研究基本过程的修改和扩展，并把它们用于游戏和单智能体应用。

\subsection*{8.12\quad 本章小结}

规划需要环境模型。
\textbf{分布模型}由可能动作对应的下一个状态和奖励的概率组成；\textbf{样本模型}则按照这些概率产生单个转移和奖励。
动态规划需要分布模型，因为它使用\textbf{期望更新}，而期望更新要对所有可能的下一个状态和奖励求期望。
另一方面，样本模型是模拟与环境交互所需的；在这种交互中，可以使用许多强化学习算法所用的那类\textbf{样本更新}。
样本模型通常比分布模型容易获得得多。

我们提出了一种观点，强调规划最优行为与学习最优行为之间出人意料地紧密的关系。
二者都涉及估计同样的价值函数；在两种情况下，自然做法都是通过一长串小的备份操作来增量地更新估计。
这使学习过程和规划过程很容易整合：只需允许二者更新同一个估计价值函数。
此外，任何学习方法都可以变成规划方法，只要把它应用于模拟的、由模型生成的经验，而不是真实经验。
在这种情况下，学习和规划甚至更加相似；它们可能是相同算法，只是作用于两个不同经验来源。

把增量规划方法与行动和模型学习整合起来也很直接。
规划、行动和模型学习以循环方式相互作用（如第 162 页的图所示），每个过程都产生另一个过程改进所需要的东西；它们之间既不需要其他相互作用，也不禁止其他相互作用。
最自然的方法是让所有过程异步并行进行。
如果这些过程必须共享计算资源，那么分配几乎可以任意处理，采用对当前任务最方便、高效的组织即可。

本章触及了状态空间规划方法之间若干变化维度。
一个维度是更新大小的变化。
更新越小，规划方法就越增量。
最小的更新包括一步样本更新，如 Dyna 中那样。
另一个重要维度是更新分布，也就是搜索聚焦的位置。
优先扫掠向后聚焦于那些价值最近发生变化的状态的前驱。
同策略轨迹采样聚焦于智能体控制环境时可能遇到的状态或状态--动作对。
这能让计算跳过与预测或控制问题无关的状态空间部分。
实时动态规划作为价值迭代的同策略轨迹采样版本，说明了这种策略相对于传统基于扫掠的策略迭代的一些优点。

规划还可以从相关状态向前聚焦，例如从智能体--环境交互中实际遇到的状态开始。
最重要的形式是决策时规划，也就是把规划作为动作选择过程的一部分。
人工智能中研究的经典启发式搜索就是一个例子。
其他例子包括 rollout 算法和蒙特卡洛树搜索；它们受益于在线的、增量的、基于样本的价值估计和策略改进。

\subsection*{8.13\quad 第一部分小结：维度}

本章结束了本书第一部分。
在这一部分中，我们试图把强化学习呈现为一组贯穿各种方法的连贯思想，而不是一堆彼此分离的方法。
每个思想都可以看作一个方法会沿其变化的维度。
这些维度的集合张成了一个很大的可能方法空间。
我们希望在维度层面探索这个空间，以获得最广泛、最持久的理解。
本节使用方法空间中维度的概念，概括到目前为止本书所发展的强化学习观点。

到目前为止，本书探索的所有方法都有三个共同关键思想：第一，它们都试图估计价值函数；第二，它们都沿着实际或可能的状态轨迹备份价值；第三，它们都遵循广义策略迭代（GPI）的一般策略，也就是说，它们维护一个近似价值函数和一个近似策略，并持续试图以彼此为基础改进二者。
这三个思想是本书所覆盖主题的核心。
我们认为，价值函数、备份式价值更新以及 GPI，是强大的组织原则，可能与任何智能模型有关，无论是人工的还是自然的。

图 8.11 显示了方法变化的两个最重要维度。
这些维度与用来改进价值函数的更新类型有关。
水平方向表示使用样本更新（基于一条样本轨迹）还是期望更新（基于可能轨迹的分布）。
期望更新需要分布模型，而样本更新只需要样本模型，或者可以完全没有模型、直接从真实经验中完成（这是另一个变化维度）。
图 8.11 的垂直方向对应于更新深度，即自举程度。
这个空间四个角中的三个，分别是估计价值的三类主要方法：动态规划、TD 和蒙特卡洛。
空间左边缘是样本更新方法，从一步 TD 更新到完整回报蒙特卡洛更新。
它们之间是一条谱系，包括基于 $n$ 步更新的方法；第 12 章还会把它扩展到 $n$ 步更新的混合，例如由资格迹实现的 $\lambda$-更新。

\begin{figure}[htbp]
\centering
\includegraphics[width=0.78\textwidth]{tex/强化学习/figures/math-9/figure-8-11.png}
\caption{强化学习方法空间的一个切片，突出本书第一部分探索的两个最重要维度：更新的深度与宽度。}
\end{figure}

动态规划方法位于空间的右上角极端，因为它们涉及一步期望更新。
右下角是极深期望更新的极端情形：它们一直运行到终止状态（或在持续性任务中，直到折扣把任何进一步奖励的贡献降到可忽略的水平）。
这就是穷尽搜索的情形。
沿这个维度的中间方法包括启发式搜索以及搜索并更新到有限深度、也许是选择性深度的相关方法。
还有一些方法在水平方向上是中间的。
它们包括混合期望更新和样本更新的方法，也包括在单次更新中混合样本与分布的可能方法。
正方形内部被填充，以表示所有这些中间方法的空间。

本书强调的第三个维度，是同策略方法与离策略方法之间的二元区分。
在前者中，智能体学习它当前遵循策略的价值函数；在后者中，它为另一个策略学习价值函数，这个策略通常是智能体当前认为最好的策略。
由于需要探索，生成行为的策略通常不同于当前认为最好的策略。
这个第三维可以想象为垂直于图 8.11 页面平面。

除刚才讨论的三个维度之外，本书还识别出若干其他维度：

\noindent\textbf{回报定义}\quad 任务是分幕的还是持续的？是折扣的还是不折扣的？

\noindent\textbf{动作价值、状态价值与后状态价值}\quad 应估计哪类价值？如果只估计状态价值，那么动作选择需要一个模型或一个单独策略（如 actor--critic 方法中那样）。

\noindent\textbf{动作选择/探索}\quad 如何选择动作，以确保在探索与利用之间取得合适权衡？我们只考虑了最简单的几种方式：$\varepsilon$-贪心、乐观初始价值、soft-max 和置信上界。

\noindent\textbf{同步与异步}\quad 所有状态的更新是同时进行，还是按某个顺序逐个进行？

\noindent\textbf{真实与模拟}\quad 应该基于真实经验还是模拟经验进行更新？如果二者都用，各用多少？

\noindent\textbf{更新位置}\quad 应更新哪些状态或状态--动作对？无模型方法只能在实际遇到的状态和状态--动作对中选择，但基于模型的方法可以任意选择。这里有许多可能性。

\noindent\textbf{更新时间}\quad 更新应作为动作选择的一部分完成，还是只在之后完成？

\noindent\textbf{更新记忆}\quad 更新后的价值应保留多久？是永久保留，还是只在计算动作选择时保留，如启发式搜索中那样？

当然，这些维度既不穷尽，也不互斥。
单个算法在许多其他方面也会不同，而且许多算法沿若干维度处在若干位置。
例如，Dyna 方法同时使用真实经验和模拟经验来影响同一个价值函数。
同样，考虑用不同方式或在不同状态与动作表示上计算的多个价值函数，也是完全合理的。
不过，这些维度确实构成了一组连贯思想，可用于描述并探索一个宽广的可能方法空间。

这里没有提到、也没有在本书第一部分覆盖的最重要维度，是函数近似。
函数近似可以看作一个正交的可能性谱系：从一个极端的表格方法，到状态聚合、多种线性方法，再到多样的非线性方法。
这个维度将在第二部分中探索。

\section*{书目与历史注释}

\noindent\textbf{8.1}
这里提出的规划与学习总体观点，是作者们多年逐渐发展出来的，部分见 Sutton（1990, 1991a, 1991b）、Barto, Bradtke, and Singh（1991, 1995）、Sutton and Pinette（1985）以及 Sutton and Barto（1981b）。
它也受到 Agre and Chapman（1990；Agre, 1988）、Bertsekas and Tsitsiklis（1989）、Singh（1993）等工作的强烈影响。
作者还深受潜伏学习心理学研究（Tolman, 1932）以及关于思维本质的心理学观点影响。
本书第三部分第 14.6 节会把基于模型和无模型方法与学习和行为的心理学理论联系起来，第 15.11 节讨论大脑可能如何实现这些方法的想法。

\noindent\textbf{8.2}
我们用来描述不同强化学习类型的 direct 和 indirect 术语来自自适应控制文献。
system identification 一词也在自适应控制中用于我们所称的模型学习。
Dyna 架构归功于 Sutton（1990），本节和下一节中的结果基于那里报告的结果。
Barto and Singh（1990）讨论了比较直接与间接强化学习方法时的一些问题。
把 Dyna 扩展到线性函数近似的早期工作由 Sutton, Szepesvári, Geramifard, and Bowling（2008）以及 Parr, Li, Taylor, Painter-Wakefield, and Littman（2008）完成。

\noindent\textbf{8.3}
若干基于模型的强化学习工作把探索奖励和乐观初始化思想推到其逻辑极限：所有未完全探索的选择都被假定为最大奖励，并计算最优路径来测试它们。
Kearns and Singh（2002）的 E$^3$ 算法和 Brafman and Tennenholtz（2003）的 R-max 算法保证能在关于状态数和动作数的多项式时间内找到近似最优解。
这通常对实际算法来说太慢，但大概是最坏情形下能做到的最好结果。

\noindent\textbf{8.4}
优先扫掠由 Moore and Atkeson（1993）以及 Peng and Williams（1993）同时独立发展。
第 170 页框中的结果归功于 Peng and Williams（1993），第 171 页框中的结果归功于 Moore and Atkeson。
这一领域后续关键工作包括 McMahan and Gordon（2005）以及 van Seijen and Sutton（2013）。

\noindent\textbf{8.5}
本节深受 Singh（1993）实验的影响。

\noindent\textbf{8.6--7}
轨迹采样从一开始就隐含地属于强化学习，但 Barto, Bradtke, and Singh（1995）在引入 RTDP 时最明确地强调了它。
他们认识到 Korf（1990）的 learning real-time A*（LRTA*）算法是一种异步 DP 算法，既适用于随机问题，也适用于 Korf 所关注的确定性问题。
除 LRTA* 外，RTDP 还包括在执行动作之间的时间间隔中更新许多状态价值的选项。
Barto 等（1995）把 Korf（1990）关于 LRTA* 的收敛证明与 Bertsekas（1982）以及 Bertsekas and Tsitsiklis（1989）保证不折扣随机最短路径问题中异步 DP 收敛的结果结合起来，证明了这里描述的收敛结果。
把模型学习与 RTDP 结合起来称为 Adaptive RTDP，也由 Barto 等（1995）提出，并由 Barto（2011）讨论。

\noindent\textbf{8.9}
关于启发式搜索的进一步阅读，建议读者参考 Russell and Norvig（2009）以及 Korf（1988）等教材和综述。
Peng and Williams（1993）探索了与本节建议相近的更新向前聚焦。

\noindent\textbf{8.10}
Abramson（1990）的 expected-outcome model 是应用于二人游戏的 rollout 算法，其中两个模拟玩家都随机行动。
他认为，即使在随机玩法下，它也是一种“强大启发式”，具有“精确、准确、易估计、高效可计算且与领域无关”的特点。
Tesauro and Galperin（1997）展示了 rollout 算法在改进双陆棋程序走法上的有效性，并采用 “rollout” 一词来表示通过用不同随机掷骰序列展开局面来评价双陆棋局面。
Bertsekas, Tsitsiklis, and Wu（1997）考察了应用于组合优化问题的 rollout 算法，Bertsekas（2013）综述了它们在离散确定性优化问题中的使用，并指出它们“常常出人意料地有效”。

\noindent\textbf{8.11}
MCTS 的核心思想由 Coulom（2006）以及 Kocsis and Szepesvári（2006）引入。
它们建立在这些作者所综述的早期蒙特卡洛规划算法研究之上。
Browne 等（2012）是关于 MCTS 方法及其应用的优秀综述。
David Silver 对本节的思想和表述有所贡献。

\end{document}
